% Generated draft: requires mathematical and source review.
\documentclass[11pt]{article}
\usepackage[margin=25mm]{geometry}
\usepackage{fontspec}
\setmainfont{Latin Modern Roman}
\usepackage{amsmath,amssymb,mathtools,unicode-math}
\setmathfont{Latin Modern Math}
\usepackage{enumitem,needspace,xurl,hyperref}
\setlength{\parindent}{0pt}
\setlength{\parskip}{5pt}
\setlength{\emergencystretch}{3em}
\newcommand{\N}{\mathbb{N}}
\newcommand{\Z}{\mathbb{Z}}
\newcommand{\Q}{\mathbb{Q}}
\newcommand{\R}{\mathbb{R}}
\newcommand{\C}{\mathbb{C}}
\newcommand{\F}{\mathbb{F}}
\newcommand{\A}{\mathbb{A}}
\newcommand{\PP}{\mathbb{P}}
\newcommand{\E}{\mathbb{E}}
\newcommand{\eps}{\varepsilon}
\newcommand{\dd}{\,\mathrm{d}}
\begin{document}
\section*{UnsolvedMath ID: 2523}
\textbf{Source problem code:} catalogue entry
\subsection{Definitions and mathematical statement}
Let $G$ be an arbitrary group with identity $1_G$. An endomorphism is a homomorphism from a group to itself. The trivial endomorphism of $G$ maps every element to $1_G$, and an inner automorphism has the form $g\mapsto tgt^{-1}$ for a fixed $t\in G$.

Suppose $\alpha\in\operatorname{End}(G)$ has the universal extension property
\[
 \forall H\text{ group}\ \forall f\in\operatorname{Hom}(G,H)\
 \exists\beta_f\in\operatorname{End}(H),\qquad
 \beta_f\circ f=f\circ\alpha.
\]
The question is whether this forces
\[
 \bigl(\forall g\in G,\ \alpha(g)=1_G\bigr)
 \quad\text{or}\quad
 \bigl(\exists t\in G\ \forall g\in G,\ \alpha(g)=tgt^{-1}\bigr).
\]
There are no finiteness, generation, injectivity, or surjectivity assumptions on the groups or on $f$. The extension may depend on $f$ and need not be unique or an automorphism; no compatibility between the different $\beta_f$ is required. The stated alternatives are not required to be disjoint in degenerate groups.

The supplied literature triage checked 17 August 2026 reports an affirmative answer from Fournier-Facio's 2026 preprint. It describes the stronger detecting-embedding theorem as follows: for every group $G$ there are a group $H$ and an embedding $i:G\hookrightarrow H$ such that a nontrivial $\gamma\in\operatorname{End}(G)$ admits some $\delta\in\operatorname{End}(H)$ with $\delta i=i\gamma$ if and only if $\gamma$ is inner. Applying the hypothesis to this particular embedding gives the requested dichotomy. The trivial map is already an allowed alternative. Conversely, a trivial $\alpha$ extends by the trivial endomorphism of every $H$, and conjugation by $t$ extends along $f$ by conjugation by $f(t)$, so both alternatives have the stated property.

The source identifies this as Question 5 of Bergman's 2012 paper and as Kourovka Notebook, New Problems, issue 21 (2026), Problem 21.14, proposed by G. M. Bergman. Both the main registry and the supplied triage label it solved. The triage explicitly qualifies the exact-match result as a recent preprint awaiting independent peer review or publication; it says no case of the stated dichotomy remains if the cited theorem is correct. This draft preserves that source qualification and makes no independent proof-verification claim.
\subsection{Short English statement}
If an endomorphism extends across every homomorphism from its group, must it be trivial or conjugation by a group element? The supplied 2026 preprint reports yes, via one embedding that detects all nontrivial inner endomorphisms, with peer review still pending.
\subsection{Sources}
\begin{itemize}
\item Local UnsolvedMath registry, permanent ID $2523$, exact source code \texttt{catalogue entry}, ``Kourovka Notebook Problem 21.14,'' including the literature triage checked 17 August 2026.
\item \emph{Kourovka Notebook}, New Problems, 21st issue (2026), Problem 21.14, attributed to G. M. Bergman in the supplied source.
\item G. M. Bergman, Publicationes Matem\`atiques \textbf{56} (2012), 91--126, Question 5, as supplied; no title or URL is given for this historical paper.
\item Francesco Fournier-Facio, \emph{An improved characterisation of inner automorphisms of groups}, arXiv:2604.05728 (2026), \url{https://arxiv.org/abs/2604.05728}; the source identifies its abstract's detecting-embedding theorem as giving the requested conclusion, with ordinary independent review/publication still pending.
\end{itemize}
\end{document}
